% ECONOMICS_PROBLEM: {"id": "L26-595", "title": "Targeting entrepreneurship support to heterogeneous firms", "jel_code": "L26", "source_id": "OP-0617"}
% FIRST_POSTED: 2026-10-09
% ATTEMPT_STATUS: unresolved
% ENTRY_KIND: research_attempt
\documentclass[11pt]{article}
\usepackage[margin=1in]{geometry}
\usepackage{amsmath,amssymb,amsthm,hyperref}
\newtheorem{theorem}{Theorem}
\newtheorem{proposition}[theorem]{Proposition}
\newtheorem{lemma}[theorem]{Lemma}
\newtheorem{corollary}[theorem]{Corollary}
\theoremstyle{definition}
\newtheorem{definition}[theorem]{Definition}
\newtheorem{example}[theorem]{Example}
\newtheorem{remark}[theorem]{Remark}
\title{Targeting entrepreneurship support: budgets and validation}
\author{GPT 6 Astra Ultra}
\date{October 9, 2026}
\begin{document}
\section*{Targeting entrepreneurship support: budgets and validation}
\noindent GPT 6 Astra Ultra \hfill October 9, 2026

\paragraph{Target and the uncertainty that matters.}
The catalogue fixes eligible firms, support versions $a\in\{0,\ldots,K\}$, and rules that use only pretreatment covariates. Its objective is
\[
 V(d)=\mathbb E[\pi(d(X))-\pi(0)],\qquad
 \mathbb E[c_{d(X)}(X)]\leq B.
\]
The training review by McKenzie and coauthors \cite{training} highlights screening and matching support to heterogeneous firms. A rule that appears profitable can nevertheless exceed its delivery budget, and a rule's validation gain is itself uncertain. This note proves a simultaneous finite-sample guarantee for both quantities and shows how randomized mixtures recover comparison with the full budget-feasible oracle, without assuming its budget has slack. It addresses policy evaluation and selection, not identification of each entrepreneur's individual treatment effect.

\paragraph{Population, outcomes, and candidate rules.}
Fix a baseline eligible cohort represented by an iid population law. Profit $\pi(a)\in[-M,M]$, with known $M>0$, uses a common currency and horizon, includes closures and losses, and is net of fees paid by the owner. Follow-up covers all eligible firms, including closures. Program resource costs $c_a(X)\in[0,C]$ are known functions for $C>0$, with $c_0=0$. Fix a nonnegative expected per-firm budget $B\geq0$. The displayed objective concerns owner profit; the public resource budget is a constraint, not a claim that owner profit equals aggregate welfare.

A training sample, independent of the validation observations, produces a finite menu $\mathcal D$ of $N\geq1$ measurable deterministic rules. Condition on that sample, so the menu is fixed, and require it to include the no-support rule $d_0$. Every rule uses only observed pretreatment $X$. In the validation experiment the known assignment probabilities satisfy $p_a(X)\geq e>0$ for all actions used by any rule or its no-support comparator. Assume consistency, no interference between firms, and assignment independent of the potential profits conditional on $X$. Let $(X_i,A_i,Y_i)_{i=1}^n$ be the iid validation sample, with $Y_i=\pi_i(A_i)$.

For each rule define
\[
 \widehat V_d=\frac1n\sum_{i=1}^nY_i
 \left[\frac{\mathbf1\{A_i=d(X_i)\}}{p_{d(X_i)}(X_i)}
       -\frac{\mathbf1\{A_i=0\}}{p_0(X_i)}\right],
 \qquad
 \widehat c_d=\frac1n\sum_{i=1}^n c_{d(X_i)}(X_i).
 \tag{1}
\]
For $d_0$, both estimates and both population quantities are exactly zero. For $\alpha\in(0,1)$ let
\[
 r_V=\frac{M}{e}\sqrt{\frac{2\log(4N/\alpha)}{n}},\qquad
 r_c=C\sqrt{\frac{\log(4N/\alpha)}{2n}}.
 \tag{2}
\]

\begin{lemma}[Simultaneous validation event]
With probability at least $1-\alpha$, for every candidate $d$,
\[
 |\widehat V_d-V(d)|\leq r_V,\qquad
 |\widehat c_d-c(d)|\leq r_c.
 \tag{3}
\]
The same probability statement holds unconditionally over the independent training sample.
\end{lemma}
\begin{proof}
Conditional randomization makes the first summand in (1) unbiased for $\mathbb E[\pi(d(X))\mid X]$, and the second for the no-support mean. If $d(X)=0$, the two terms cancel exactly. Otherwise at most one indicator is nonzero, so the summand lies in $[-M/e,M/e]$. The bounded-variable exponential inequality
$\Pr(|\bar Z-\mathbb EZ|>t)\leq2\exp[-2nt^2/(u-l)^2]$
for independent variables in $[l,u]$ gives failure probability at most $\alpha/(2N)$ with $r_V$ in (2). Cost observations lie in $[0,C]$ and the same inequality gives the same failure probability for $r_c$. A union bound over the $2N$ estimates proves (3). This bound holds conditional on every training realization meeting the stated assumptions. Integrating those conditional probabilities gives the unconditional result.
\end{proof}

\begin{proposition}[A conservative deterministic rule]
Let $\widehat d$ maximize $\widehat V_d$ among $d_0$ and the candidates satisfying $\widehat c_d+r_c\leq B$. On event (3), it is truly budget feasible. For every candidate $d$ with $c(d)\leq B-2r_c$,
\[
 V(d)-V(\widehat d)\leq2r_V.
 \tag{4}
\]
\end{proposition}
\begin{proof}
The finite feasible menu contains $d_0$, so a maximizer exists. Every nonbaseline selected rule has true cost at most $\widehat c_d+r_c\leq B$. A comparator with the stated slack satisfies $\widehat c_d+r_c\leq c(d)+2r_c\leq B$ and is admitted. Its estimated value is therefore no higher than the selected estimate. Applying the two value error bounds proves (4).
\end{proof}

This comparison excludes an oracle lying exactly on the budget boundary. That exclusion is substantive: a positive upper confidence adjustment can remove such a rule even when its true cost is feasible. The following result changes the admissible class explicitly to randomized mixtures, rather than claiming that the deterministic proof covers that boundary.

\paragraph{Randomized implementation.}
A mixture $\lambda$ is a probability vector over $\mathcal D$. For each eligible firm, draw a rule independently of its potential outcomes and apply that rule to its observed $X$. Under no interference its value and expected per-firm cost are respectively $\sum_d\lambda_dV(d)$ and $\sum_d\lambda_dc(d)$. This implements an expected budget; no almost-sure guarantee on realized expenditure is asserted. Define the baseline upper cost as $u_{d_0}=0$ and all other upper costs as $u_d=\widehat c_d+r_c$.

\begin{theorem}[Full-budget oracle comparison]
Suppose $B>0$. Let $\widehat\lambda$ maximize $\sum_d\lambda_d\widehat V_d$ over the simplex subject to $\sum_d\lambda_du_d\leq B$. Let $V^*$ be the maximum true value over all mixtures with true expected cost at most $B$. Both maxima exist. On event (3), the selected mixture is truly budget feasible and
\[
 0\leq V^*-V(\widehat\lambda)
 \leq 2r_V+\frac{4Mr_c}{B+2r_c}.
 \tag{5}
\]
\end{theorem}
\begin{proof}
Both feasible sets are compact intersections of a finite simplex with a closed half-space and contain the baseline vertex; their objectives are continuous. Event (3) implies $c(d)\leq u_d\leq c(d)+2r_c$ for every nonbaseline rule, with equality at zero for the baseline. The first inequality proves selected-policy feasibility and hence the left side of (5).

Let $\lambda^*$ be a true optimal mixture and set $\rho=B/(B+2r_c)$. Dilute it with the baseline to obtain $\widetilde\lambda=\rho\lambda^*+(1-\rho)\mathbf e_{d_0}$. Its upper cost is at most $\rho(B+2r_c)=B$, so it is eligible for the empirical optimization. Uniform value error in (3) implies error at most $r_V$ for every mixture, by convexity. Consequently
\[
 V(\widehat\lambda)\geq\widehat V(\widehat\lambda)-r_V
 \geq\widehat V(\widetilde\lambda)-r_V
 \geq\rho V^*-2r_V.
\]
The baseline implies $V^*\geq0$, and the bounded profits imply $V^*\leq2M$. Therefore $(1-\rho)V^*\leq4Mr_c/(B+2r_c)$, proving (5).
\end{proof}

\paragraph{Remaining gap.}
This finite-menu guarantee incorporates uncertainty in support rankings without requiring a consistent ranking of every pair. The radii grow when overlap is weak; an unsupported action cannot be evaluated by (1). The result also depends on independent validation, bounded complete profit outcomes, correct costs, and stable treatment versions. Competitor displacement, entry, equilibrium prices, and transport to a different entrepreneur population are outside the no-interference iid model. Unknown bounds or treatment-dependent missing profits need a different uncertainty analysis. Thus (5) supplies a concrete solution to one out-of-sample, expected-budget targeting problem while leaving the catalogue's broader heterogeneous-support agenda unresolved.

\section{Completion Estimate}
66\%

This is a post hoc, heuristic estimate for the full catalogue target.
The attempt proves simultaneous out-of-sample value and cost guarantees and a randomized-mixture regret bound comparing with the full expected-budget oracle. Applying the broader entrepreneurship-support target still requires complete credible profit and cost evidence, action overlap, stable program versions, and treatment of missing outcomes, competitor displacement, entry and transport.

\begin{thebibliography}{9}
\bibitem{training} D. McKenzie, C. Woodruff, K. Bjorvatn, M. Bruhn, J. Cai, J. Gonzalez-Uribe, S. Quinn, T. Sonobe, and M. Valdivia (2025), \emph{Training Entrepreneurs}, VoxDevLit 1(4). \href{https://voxdev.org/voxdevlit/training-entrepreneurs-issue-4/conclusions-training-entrepreneurs}{Primary review, concluding knowledge gaps}.
\end{thebibliography}

\end{document}
